\documentclass[11pt,aspectratio=169]{beamer}

% --------------------
% Theme & Packages
% --------------------
% \usetheme{Boadilla}
\usetheme{Madrid}
\usecolortheme{default}
\usefonttheme{professionalfonts}  % 数学公式字体
% \setbeamercolor{block title}{fg = white,bg = blue}
\usepackage{ctex}
\usepackage{amsmath, amssymb}
\usepackage{listings}
\usepackage{xcolor}

% --------------------
% Code Style
% --------------------
\lstset{
    language=C++,
    basicstyle=\ttfamily\small,
    keywordstyle=\color{blue},
    stringstyle=\color{red},
    commentstyle=\color{gray},
    numbers=left,
    numberstyle=\tiny,
    stepnumber=1,
    breaklines=true,
    frame=single
}

% --------------------
% Title Info
% --------------------
\title{导论}
\author{qwaszx}
\date{\today}

% ====================
% Document
% ====================
\begin{document}

% --------------------
\begin{frame}
  \titlepage
\end{frame}

% --------------------
\begin{frame}{今天的内容}
\begin{itemize}
    \item 为什么学 OI / 学 OI 有什么好处？
    \item 怎么学 OI？
    \item 回顾一些语法和算法基础
    \item 对本系列课程的介绍
\end{itemize}
\end{frame}

% =========================================================
% Part 1
% =========================================================
\section{为什么要学 OI}

\begin{frame}{什么是 OI？}

\pause
面向中学生的\textbf{算法竞赛}：设计算法解决问题

通常来说，需要使用\textbf{代码}实现算法，最后使用若干组数据进行测试。

\pause

\begin{exampleblock}{不是}
    \begin{itemize}
        \item 写代码竞赛
            \begin{itemize}
                \item 不过也可以是（工程题/大模拟）
            \end{itemize}
        \item 数学竞赛
            \begin{itemize}
                \item 不需要证明
                \item 只需要通过测试
            \end{itemize}
    \end{itemize}
\end{exampleblock}
\end{frame}

\begin{frame}{为什么学了 OI}
    你为什么学了 OI？

    \pause
    我为什么学了 OI。
\end{frame}

\begin{frame}{学 OI 的好处}
\begin{itemize}
    \item 提升逻辑思维能力
    \item 算法思维
    \begin{itemize}
        \item 如何描述解决问题的步骤
        \item 对计算的理解
    \end{itemize}
    \item 计算机和离散数学知识
    \item OI 社群
    \item 更多地玩电脑
    \item 升学与就业
\end{itemize}
\end{frame}

\begin{frame}{升学与就业}
升学方面：
    \begin{itemize}
        \item 金牌可以保送清华姚班/北大图灵班
        \item 银牌可以破格入围强基计划
        \item 省一对强基计划和各种自主招生有一定帮助 
    \end{itemize}
就业方面：
    \begin{itemize}
        \item OI/XCPC 竞赛获奖对于计算机相关行业有一定帮助
        \item 可以继续做教培，相对普通文化课报酬更高
        \item 对口 TCS（理论计算机科学）
    \end{itemize}
\end{frame}

\begin{frame}{好处说完了，那么坏处呢}
    \begin{itemize}
        \item 机会成本：选择 OI 需要放弃一些其他东西
        \item 更多地使用计算机和互联网，可能会损害身心健康。
        \item 对于升学和就业的性价比不算很高
        \begin{itemize}
            \item 目前，银牌相比省一的优势不大，但是困难得多
            \item 集训队每年只有 50 人，竞争激烈
        \end{itemize}
    \end{itemize}
\end{frame}

\begin{frame}{总结}
    初中阶段，适合学文化课有余力并且学 OI 学得开心的同学学习。

    高中阶段，功利角度建议根据自身成绩谨慎选择（高一进省队 -> 高二进集训队）

    兴趣是第一位的！
\end{frame}
% =========================================================
% Part 2
% =========================================================
\section{怎么学 OI}

\begin{frame}{OI 比赛的级别}
    \begin{enumerate}
        \item CSP-J
        \item CSP-S
        \item NOIP
        \item 省选
        \item NOI
        \item CTS
        \item IOI
    \end{enumerate}
    其中 CSP-J ~ CSP-S 属于前期，CSP-S ~ NOIP 属于中期，省选 ~ NOI 属于后期，再往后是大后期
\end{frame}

\begin{frame}{做题的过程}
\begin{enumerate}
    \item \textbf{读题}
    \item \textbf{分析性质}
    \item \textbf{抽象问题}
    \item \textbf{使用/改编/设计算法来处理}
    \item \textbf{实现}
\end{enumerate}

前期最重要的是 1 和 5，中期最重要的是 3 和 4，后期和大后期最重要的是 2 和 4

学习的过程中，每个部分都在逐步提高
\end{frame}

\begin{frame}{读题与实现}
    抽象/封装：忽略某个东西的内部结构，只关注其输入和输出

    在做题中练习

    多看看其他人的写法
\end{frame}

\begin{frame}{抽象问题，调用算法}
    我们会学习很多的算法。学习的过程中，
    \begin{itemize}
        \item 关注算法整体的功能，熟知一个问题可以被什么算法解决，一个算法可以解决什么问题，有哪些变种和扩展
        \item 关注算法的每一步，熟知其原理、作用，多思考可以如何扩展
        \item 关注证明
    \end{itemize}

    在做题中练习

    多总结
\end{frame}

\begin{frame}{分析性质，设计算法}
    也就是所谓的人类智慧。一般来说，很难！

    多总结题目和算法的模型。相似的模型往往具有相似的性质。

    多思考题解/算法的每一步是如何想到的。

    尝试建立不同概念和不同算法之间的联系，抽象其共性
\end{frame}

\begin{frame}{训练与学习资源}
    OJ/题库：洛谷、QOJ、LOJ、UOJ；Codeforces、Atcoder

    比赛：Codeforces、Atcoder

    学习资源：oiwiki、各种优质博客（如 liu\_cheng\_ao 推荐的、command\_block、Elegia）

    各种 AI

    OI 社群与搜索引擎。不懂就问/搜！
\end{frame}

\begin{frame}{训练的方法}
    \begin{itemize}
        \item 学习知识点
        \item 刷知识点对应的题
        \item 对知识点进行归纳总结
        \item 通过模拟赛进行不同知识点的综合练习/复习
    \end{itemize}
\end{frame}

\begin{frame}{刷题的方法}
    先自行思考，思考到自己一直没有任何想法为止，然后看题解。这个时间长短是不定的

    看题解的过程中也要进行思考：作者说的有没有道理，作者是怎么想到的，我能不能自行补全剩下的部分，哪些部分还可以优化

    想出做法/看懂题解后进行实现，这个过程中不要参考任何现成代码

    做完题后看看题解区有没有不同做法，看看其他人的代码都是怎么写的，想想自己的代码哪些地方能更快/更好写

    看看题目是否是\textbf{经典问题}，是则需要在归纳总结时\textbf{记录}
\end{frame}

\begin{frame}{经典问题}
    有一些问题是经典的：在拆掉包装后，其技巧/模型非常泛用。每做完一道题都应该思考这道题是不是这种问题。换句话说，做题后要努力将题目的模型/技巧总结出来。
\end{frame}

\begin{frame}{刷什么题}
    学习资源中提到的题、自行根据知识点搜索相关题目（如洛谷题单）

    重要的是刷题前中后的\textbf{思考}而不是刷题本身。所以，非常同质化的题不用多刷，几道就够了（再多的可以口胡）

    刷题可以形成题目<->模型<->算法的条件反射，或者说直觉。
\end{frame}

\begin{frame}{训练的理念}
    循序渐进，避免一直学太简单/太难的东西

    青少年的智力、理解能力会随着年龄提高

    熟练掌握 OI 大部分知识点通常需要且只需要几千个小时，相当于 1~2 年*8h（高中脱产/初中或者小学就开始学） 
\end{frame}

\begin{frame}{OI 之外}
    对初中生来说，我非常建议学好文化课，义务教育是必要的。不过可以多和家长老师沟通，放弃重复性过强的文化课练习。OI 以兴趣为主，无需强求。

    OI 涉及大量的思考，所以充足睡眠非常重要，建议想办法得到充足睡眠。最起码的原则是不能一天比一天困。

    多锻炼。一方面大学也有各种体育测试（如 3km），另一方面人生是场马拉松，健康地活着才是成功

    人际交流，线上线下均可
\end{frame}

% =========================================================
% Part 3
% =========================================================
\section{语法基础 Revisit}

\begin{frame}{代码能力}
    可以用代码实现想法

    知道一个东西能不能写、大概是怎么写的、好不好写
    
    不出 bug 地尽量快地写出来

    学习更优秀的写法（更好写/更快）
\end{frame}

\begin{frame}{先想再写 vs 边写边想}
    我的习惯：在写某个部分时，对这个部分的整体结构必须提前想好，细节可以边写边想

    抽象/封装的代码设计理念
\end{frame}

\begin{frame}{目前你需要非常熟练的东西}
    \begin{itemize}
        \item 顺序、分支、循环结构的程序设计
        \item 各种变量类型、结构体、数组
        \item 函数
        \item 简单的模拟、枚举
    \end{itemize}

    如果不熟练，可以先多做做橙题
\end{frame}

% =========================================================
% Part 4
% =========================================================
\section{算法基础 Revisit}

\begin{frame}{什么是算法}
    解决问题的\textbf{机械}流程。任何人/机器只要照着做就能成功。

    流程如何描述取决于执行者：C++/python 等各种编程语言；人类；图灵机
\end{frame}

\begin{frame}{复杂度}
    流程需要的\textbf{基本指令}数量，其中每个基本指令花费 1 个单位时间。

    什么是\textbf{基本指令}？
    \begin{itemize}
        \item word RAM——通常的选择：算术运算、位运算、内存访问
        \item （不严谨）人类的一个动作
        \item 图灵机移动一格（想象数组不能直接访问，而是只能访问上一个或者下一个位置）
    \end{itemize}

    现代计算机中，不同指令的速度很不相同，这对算法的实际运行效率有很大影响。可以询问 ai “现代计算机的各种操作的速度如何排序”。目前 CCF 的机器 1s 大概可以跑 $10^9$ 个时钟周期。
\end{frame}

\begin{frame}{大 $O$ 记号}
    $g(n)$ 是 $O(f(n))$ 意味着：存在 $N$ 和 $C>0$，使得对任何 $n>N$，都有 $g(n)\le C\cdot f(n)$。换句话说，当 $n$ 非常大时，$g(n)$ 不会超过常数倍的 $f(n)$。
    
    讨论算法复杂度时，我们通常使用渐进记号来描述，其中 $n$ 代表问题规模。直观来说，我们只考虑 $f(n)$ 中关于 $n$ 增长最快的一项，且忽略其常数。

    OI 中常见的复杂度：$O(1)$、$O(\log n)$、$O(\sqrt{n})$、$O(n)$、$O(n\log n)$、$O(n\log^c n)$、$O(n^c)$、$O(2^n n^c)$、$O(c^n n^k)$
\end{frame}
\begin{frame}{渐进记号的各种用法}
    $2^{O(\log n)}$ 是什么意思？

    人们在表达式的某个位置中写 $O(f(n))$ 时，表示：存在某个常数 $C>0$，使得这里的量在 $n$ 非常大的时候不超过 $C\cdot f(n)$。

    带常数的复杂度：$2\log n+O(1)$

    数学中的例子：$\frac{n}{n+1}=1+O(1/n)$。（这里是对绝对值用 $O$）

\end{frame}
\begin{frame}{其他渐进记号}
    $\Omega(f(n))$：把大 $O$ 的 $\le C\cdot f(n)$ 换成 $\ge C\cdot f(n)$

    $o(f(n))$：当 $n$ 趋于无穷时，与 $f(n)$ 的比值趋于 $0$，也就是严格比常数倍 $f(n)$ 小

    $\omega(f(n))$：把小 $o$ 的趋于 $0$ 换成趋于 $\infty$
    
\end{frame}

\begin{frame}{枚举}
    TODO：找几个例子

    后面我们会学习搜索，来实现最一般的枚举
\end{frame}

\begin{frame}{递归}
    算法可以调用自己。在设计算法时，可以寻找“子问题”，然后调用递归

    TODO：找几个例子

    抽象贯穿计算机科学的始终。我们后面还会不断地见到子问题的思想。
\end{frame}

\begin{frame}{常见基础算法}
\begin{itemize}
    \item 二分
    \item 排序
    \item 贪心
\end{itemize}
TODO：举几个例子
\end{frame}

\section{对本课程的介绍}
\begin{frame}{培养目标}
    在课程开始前，我们希望大家具有基本的 C++ 和算法知识，可以熟练切掉橙题。

    在上半段课程中，我们会学习各种 CSP-S 级别的算法，期望帮助大家完成普转提的过渡

    在下半段课程中，我们会更深入地学习这些算法相关的内容，期望大家在这部分内容上达到 NOIP/省选/NOI 的水平（注意：这三个东西的难度其实没有显著区别）。不过由于这部分内容密度较大，学习起来可能较为困难。
\end{frame}

\begin{frame}{课程节奏}
    每周一次课 4h + 每周题单 + 一次模拟赛

    课下希望大家多多练习，大概每天做两道题左右即可
\end{frame}
% =========================================================
% Part 5
% =========================================================
\section{总结}

\begin{frame}{总结}
\begin{itemize}
    \item 对 OI 概况、利弊的了解
    \item OI 的一些学习方法
    \item 一些语法和算法基础的回顾
\end{itemize}
\end{frame}

\begin{frame}{最后}
希望大家学得开心

\pause
\centerline{\Large 谢谢大家！}
\end{frame}

\end{document}
